\section{Domain-specific languages}
\label{sec:dsl}

\begin{definition}
A \textbf{domain-specific language} (DSL) is a language designed to express
programs of one domain, in the vocabulary of that domain. A general-purpose
language (GPL) is designed for programs of any domain.
\end{definition}

Kaplan calls them little languages~\cite{kaplan1994little}. A DSL gives up
generality: its programs say what a domain expert would say, and the
implementation deals with how to do it. Most DSLs are not Turing complete,
and that restriction makes them easier to analyze and validate than programs
of a GPL. The regular languages lecture notes that regular languages are
usually used to describe finite languages such as DSLs. The course has met
several DSLs already, and this section starts from them.

\subsection{The DSLs of the course}

\paragraph{Dates.} The first DSL of the course recognizes dates in three
formats, \texttt{MM/DD/YYYY}, \texttt{DD-MM-YY} and \texttt{YYYY.MM.DD}, to
validate the values of a \texttt{Date} type. Its grammar is regular, and it is
written in Rascal with \texttt{lexical} rules (formal languages notes, section
on regular languages). The grammar checks the shape of a date but not the
calendar: day 31 of February is accepted, and rejecting it is a job for the
semantics of the DSL.

\paragraph{LDBL.} The second DSL manages the database of an image
marketplace, with the commands \texttt{ADD}, \texttt{REMOVE}, \texttt{SEARCH},
\texttt{VIEW} and \texttt{HALT}, assignment of image variables, and the
statements \texttt{if \dots{} then \dots{} else \dots{} fi} and
\texttt{loop \dots{} do \dots{} od}. The nesting of statements makes it
context-free and not regular (formal languages notes, section on context-free
languages).

\paragraph{SCuLPT.} The Simple Cubic Language for Programming Tasks, from the
first lecture, replaces typed text by tangible 3D pieces that are composed to
build programs and art pieces. Unlike the two previous languages it is Turing
complete. Its domain is not a kind of data but a kind of user and interface:
the language changes the medium in which programs are written.

\paragraph{Objectilang.} The language of the course project is an
object-oriented language with generics. A program is made of objects, which
can hold values (\texttt{val}), variables (\texttt{var}) and methods, and a
\texttt{main} method. Objects extend other objects with \texttt{extends}, and
generic objects declare each type parameter covariant with \texttt{+} or
contravariant with \texttt{-}, as in \texttt{List[+A]}. A method is attached to
an object with the \texttt{method} abstraction, as in
\texttt{Cat.talk = method () : String \{ \dots{} \}}, and \texttt{override}
replaces a method inside the construction of an object. The reserved words are
\texttt{object}, \texttt{extends}, \texttt{method}, \texttt{val}, \texttt{var},
\texttt{new} and \texttt{override}, and the control flow has
\texttt{if-then-else} and \texttt{for}, which iterates over a range or a
collection. Objectilang is closer to a GPL than the other three, but building
it touches every topic of the course: the reserved words belong to the lexer,
the nesting of objects and methods to the parser, \texttt{extends} and the
variance annotations to the type system (the subtyping section of the types
notes uses its \texttt{Transformer[-In, +Out]} example), and method lookup and
\texttt{override} to the run-time representation of objects.

\subsection{Other examples}

The following examples go beyond the slides.

\paragraph{SQL.} The domain is relational data. A query says which rows are
wanted, not how to find them:
\codelines{SQL}{sql/workouts.sql}{5}{8}
There are no loops and no iteration order. Whether the join uses a hash table,
nested loops or an index is decided by the database engine, and SQL engines
can reorder joins and filters because the queries have no side effects.

\paragraph{Make.} A Makefile describes a dependency graph: each target lists
the files it depends on and the command that builds it. The program is the
graph, not a sequence of commands; \texttt{make} orders the targets
topologically and rebuilds only those whose dependencies changed.

\paragraph{Regular expressions.} The domain is string patterns. The pattern
\codelines{Python}{python/internal.py}{28}{28}
matches a capitalized word followed by four digits, which would take a loop
and several conditions in a GPL. The density that makes regular expressions
short also makes them hard to read for someone who does not know the
notation.

\subsection{A running example: a training DSL}

The rest of these notes use a DSL for training plans, an example that goes
beyond the slides. A coach writes
\begin{lstlisting}
Monday:
  run: 10km @threshold
  rest: 15min
Wednesday:
  swim: 2000m @steady
  strength: 30min
\end{lstlisting}
and the program contains nothing that is not about training. The grammar, in
the BNF notation of the course, is
\begin{lstlisting}
Plan      ::= Day+
Day       ::= DayName ":" Activity+
Activity  ::= Exercise ":" Quantity Intensity?
Quantity  ::= Number Unit
Intensity ::= "@" Label
DayName   ::= "Monday" | "Tuesday" | ... | "Sunday"
Exercise  ::= "run" | "swim" | "cycle" | "strength" | "rest"
Unit      ::= "km" | "m" | "min" | "h"
\end{lstlisting}
Each rule can be chosen by its first token, so the grammar is LL(1). In Rascal:
\codelines{Rascal}{rascal/src/Training.rsc}{6}{19}

\subsection{A formal definition}

A DSL is described by the same three components as any language:
\[
  \mathit{DSL} = \langle \Sigma, G, S \rangle,
\]
where $\Sigma$ is the alphabet, $G$ a grammar over $\Sigma$ that defines the
valid programs, and $S$ a semantic function from parse trees (or ASTs) to
meanings. A program is a word $p \in L(G)$, and its meaning is
$S(\mathit{parse}(p))$. What changes from a GPL is the size: $\Sigma$ is small,
$G$ is simple, and $S$ only has to give meaning to the concepts of the domain.

\subsection{External and internal DSLs}

An \textbf{external} DSL has its own syntax and its own implementation: lexer,
parser, checker and interpreter or compiler. SQL, Make, the date DSL, LDBL,
Objectilang and the training DSL are external. The syntax can follow the
domain exactly, and programs can be read by people who do not program; the
whole implementation has to be built.

An \textbf{internal} (embedded) DSL is written inside a host language, with the
host's functions, method chaining or macros, and every program is also a
program of the host:
\codelines{Python}{python/internal.py}{3}{24}
The parser, type checker and debugger of the host come for free. In exchange
the program still looks like Python, and its errors are Python errors. Lisp
and Scheme are the classic hosts for internal DSLs, because their code is a
list and a macro can transform it before it runs (homoiconicity, in the
runtime systems notes). Table~\ref{tab:ext-int} compares the two.

\begin{table}[h]
\centering
\begin{tabular}{@{}p{3.2cm}p{5cm}p{5cm}@{}}
\toprule
& External & Internal \\
\midrule
Syntax & free, follows the domain & limited by the host \\
Implementation & lexer, parser, checker, evaluator & a library \\
Tools & built from scratch & those of the host \\
Error messages & in terms of the domain & in terms of the host \\
Typical users & domain experts & programmers \\
\bottomrule
\end{tabular}
\caption{External and internal DSLs.}
\label{tab:ext-int}
\end{table}

\subsection{Processing a DSL program}

An external DSL program goes through the phases studied in the course
(Figure~\ref{fig:pipeline}). The lexer groups characters into tokens:
\texttt{run: 10km @threshold} becomes the tokens \texttt{run}, \texttt{:},
\texttt{10}, \texttt{km}, \texttt{@} and \texttt{threshold}, and a character
outside the alphabet is a lexical error. The parser builds the tree according
to the grammar, and a sequence of tokens that the grammar does not generate is
a syntax error. The semantic analysis checks what the grammar cannot express,
such as the unit that each exercise accepts. The evaluator walks the checked
tree and produces the result. Each kind of error is found in a different
phase and deserves a different message.

\begin{figure}[h]
\centering
\begin{tikzpicture}[st/.style={draw, minimum width=2.1cm, minimum height=8mm, font=\small, align=center}, er/.style={draw, dashed, font=\footnotesize, align=center}, node distance=6mm and 10mm]
  \node[st] (src) {source\\ text};
  \node[st, right=of src] (lex) {lexer};
  \node[st, right=of lex] (par) {parser};
  \node[st, right=of par] (sem) {semantic\\ analysis};
  \node[st, right=of sem] (ev) {evaluator};
  \draw[flecha] (src) -- (lex);
  \draw[flecha] (lex) -- node[above, font=\scriptsize] {tokens} (par);
  \draw[flecha] (par) -- node[above, font=\scriptsize] {AST} (sem);
  \draw[flecha] (sem) -- (ev);
  \node[er, below=of lex] (e1) {lexical\\ error};
  \node[er, below=of par] (e2) {syntax\\ error};
  \node[er, below=of sem] (e3) {semantic\\ error};
  \draw[flecha, dashed] (lex) -- (e1);
  \draw[flecha, dashed] (par) -- (e2);
  \draw[flecha, dashed] (sem) -- (e3);
  \node[font=\scriptsize, above=1mm of lex] {lexical analysis};
  \node[font=\scriptsize, above=1mm of par] {parsing};
  \node[font=\scriptsize, above=1mm of sem] {types, scopes};
  \node[font=\scriptsize, above=1mm of ev] {semantics, runtime};
\end{tikzpicture}
\caption{Phases of the processing of a DSL program, labelled with the topics of the course that study them.}
\label{fig:pipeline}
\end{figure}

\subsection{Syntax, semantics and execution}

A language is defined by its syntax, which says which programs are valid; its
semantics, which says what they mean; and its execution model, which says how
the meaning becomes action. The three styles of semantics of the course apply
to a DSL as to any language. An operational rule for an activity adds it to
the state $\sigma$ of the schedule:
\[
  \langle \texttt{run: } n\,u,\ \sigma \rangle \to
  \sigma[\mathit{activities} := \mathit{activities} \cdot \mathit{Activity}(\mathit{run}, n, u)].
\]
A denotational definition maps the activity directly to a value:
\[
  \llbracket \texttt{run: 10km @threshold} \rrbracket =
  \mathit{Activity}(\mathit{run}, 10, \mathit{km}, \mathit{threshold}).
\]
An axiomatic definition would state properties of plans with pre- and
postconditions. The execution model can be an interpreter, which walks the
AST; a compiler, which translates it to lower-level code; or a translator to
another high-level language.

\subsection{Reasons and costs}

A DSL is worth building when its users are not programmers, as the coaches
of the training DSL; when the domain has rules that the language can enforce
before a program runs, as rejecting \texttt{run: 10min}; when programs should
be data that can be stored, compared and transformed; and when a restricted
language allows analyses that a GPL does not, as the optimization of SQL
queries.

The costs are the whole implementation (lexer, parser, checker, evaluator,
error messages), the lack of tools such as debuggers and editors unless they
are built too, the evolution of the language without breaking old programs,
and the time users need to learn the notation. When the domain is unstable or
the users are programmers, a library in an existing language is often enough.

\subsection{Exercises}

\begin{exercise}
Classify the date DSL, LDBL and the training DSL as regular or context-free.
\end{exercise}

\paragraph{Solution.} The date DSL is regular: every rule is a concatenation of
finite sets of symbols. LDBL is context-free and not regular, because
\texttt{if} and \texttt{loop} contain commands that can contain other
\texttt{if}s and \texttt{loop}s. The training DSL has no nesting: a plan is a
sequence of days and a day a sequence of activities, so the language is
regular, even though it is written with a context-free grammar.

\begin{exercise}
In Objectilang, \texttt{object}, \texttt{extends} and \texttt{override} are
reserved words. Which phase rejects a program that uses \texttt{extends} as the
name of a value?
\end{exercise}

\paragraph{Solution.} The lexical definition of identifiers excludes the
reserved words, so \texttt{extends} is never an identifier, and the parser
finds the keyword where the grammar expects a name: it is a syntax error. The
LDBL grammar of the formal languages notes shows the mechanism in Rascal, with
a \texttt{keyword} declaration and a reject in the lexical of names.

\begin{exercise}
Is the training DSL external or internal? Write the plan
\texttt{Monday: run: 10km @threshold rest: 15min} with the internal version in
Python.
\end{exercise}

\paragraph{Solution.} The training DSL is external: it has its own grammar and its
own parser. The internal version is
\texttt{TrainingPlan().on('Monday').run(10, intensity='threshold').rest(15)},
which is the program tested in \texttt{internal.py}.
